\section{Method: SCORE as a world compiler}\label{method-score-as-a-world-compiler}

\subsection{Overview}\label{overview}

A compiler turns a source a person wrote into an output a machine runs, through stages whose intermediate forms can be inspected, cached and checked. SCORE has the same shape. Its source is the creator's picks; its output is a world an engine loads; its stages are listed in Table 1. A world is compiled once, offline, in batches. Nothing is generated while the world is played, and each scale of a world (a room, a planet's surface, an orbit) is its own kind of world rather than one continuous zoom.

\textbf{Table 1.} The stages of the route as run for world 1. ``Who'' says what decides or does the work: the creator (a human pick), a coding agent, a model in a cloud batch, or a deterministic tool.

{\def\LTcaptype{none} % do not increment counter
\begin{longtable}[]{@{}
  >{\raggedright\arraybackslash}p{(\linewidth - 6\tabcolsep) * \real{0.2500}}
  >{\raggedright\arraybackslash}p{(\linewidth - 6\tabcolsep) * \real{0.2500}}
  >{\raggedright\arraybackslash}p{(\linewidth - 6\tabcolsep) * \real{0.2500}}
  >{\raggedright\arraybackslash}p{(\linewidth - 6\tabcolsep) * \real{0.2500}}@{}}
\toprule\noalign{}
\begin{minipage}[b]{\linewidth}\raggedright
Stage
\end{minipage} & \begin{minipage}[b]{\linewidth}\raggedright
Input
\end{minipage} & \begin{minipage}[b]{\linewidth}\raggedright
Output
\end{minipage} & \begin{minipage}[b]{\linewidth}\raggedright
Who
\end{minipage} \\
\midrule\noalign{}
\endhead
\bottomrule\noalign{}
\endlastfoot
1. References and concept & the creator's brief and reference pictures & a few concept pictures; one is picked & picture model draws; creator picks \\
2. Dimensioned plan & the picked concept & a plan in metres: shell, doorways, walkways, floor levels & coding agent; checks \\
3. Inventory & plan and concept & every object as a row: kind, anchor, size, count; parent and child rows & coding agent; concept-density check \\
4. Close-ups & inventory rows & one clean picture of each object, front-on & picture model \\
5. Code or pipeline & close-up and a code build & each kind routed to code or to the prop pipeline & parts check (deterministic) \\
6. Shape & routed kinds & code builders; or picture-to-3D models made solid & coding agent writes builders; Pixal3D in cloud batches \\
7. Surfaces & shapes with named parts & baked maps from the surface library, one room wear & ProcFunc recipes in cloud batches \\
8. Sound & surfaces, rooms, inventory & footstep, impact and echo per surface and room, a sound per object & MOSS-SoundEffect in a cloud batch; loudness rules \\
9. Assembly & everything above & the scene package an engine loads & deterministic tools; checks \\
10. Review & shots from fixed cameras & keep, or send back one stage & creator \\
\end{longtable}
}

\subsection{The canonical scene: a generated base layer and edit layers}\label{the-canonical-scene-a-generated-base-layer-and-edit-layers}

A compiler needs one intermediate form every stage agrees on. For SCORE that form is a scene description with a strict schema: every object has a name, a kind, a place and turn in a stated frame, a size in metres, a collision shape, the surface of each of its parts, its sound, and tags that say what a person can do with it (a door, a seat, a terminal, an airlock). Geometry is glTF 2.0 \citep{gltf2}. The scene itself is designed to be an OpenUSD stage \citep{openusd} in two kinds of layer: a generated base layer that the compiler owns and may rewrite, and edit layers that hold the creator's own changes. Because OpenUSD composes layers, a regeneration replaces the base layer and the creator's edits survive on top of it. Engines and tools (Godot, Blender, Unreal, a robotics simulator) then load the same scene through thin adapters.

That is the design. What runs today is narrower, and we say so plainly: in world 1 the canonical form is glTF models plus one JSON layout file per place, read by one adapter, the Godot engine's \citep{godot}. The OpenUSD stage, its layer split and every adapter other than Godot's are \textcolor{red}{[TODO: not built; the move is the next piece of work in the framework repository]}.

\subsection{Checks before spending}\label{checks-before-spending}

Every fault costs more the later it is found. A fault costs nothing in the plan, cents at the picture, about €0.18 and 75 minutes at the 3D model, and an evening of the creator's time once a room is built and played. So each paid stage is preceded by checks that are cheap, deterministic and run on real geometry rather than on an agent's judgement, and a failed check sends the work back one stage, never forward with a flag. The checks used in world 1 were:

\begin{itemize}
\tightlist
\item
  \textbf{Room checks} on the laid-out room's real meshes, about two seconds a room: rays that leave through walls or roof (leak), pieces in front of a real opening, pieces that do not lie on the surface they belong to, and the room's piece, triangle and lamp budget.
\item
  \textbf{Door check:} every door must stand in a wall or partition that fully separates its two sides, read as a walk on a grid.
\item
  \textbf{Concept-density check:} every element visible in the picked concept, cut into crops, has an inventory row or a written reason it was dropped; after the build the room is shot from the concept's own camera and compared.
\item
  \textbf{Parts check:} a code builder may only stand in for a picture-to-3D model if its build shows every part the object's clean close-up shows (Section 3.4).
\item
  \textbf{Model check:} every model is closed and solid with walls of at least 3 mm; a generated piece that should be square fails when its sides tilt more than 2 degrees or warp (the straightness check).
\item
  \textbf{Palette sweep and storage budget:} every made model's colours against the room's own library palette, and a budget per room of 250 MB on disk and 640 MB of textures on the graphics card.
\item
  \textbf{Scene check and made-only check} in the engine after install: rays that find what floats, sinks or overlaps, and a walk of the loaded scene that fails on any visible mesh the route did not make.
\end{itemize}

\subsection{Code or prop pipeline: the sorter and the parts check}\label{code-or-prop-pipeline-the-sorter-and-the-parts-check}

Picture-to-3D models are good at chunky solids and bad at flat panels, thin beams and things with openings: in the base's central room they made panels about as deep as they were wide, beams 3 to 10 times too thick, and holes filled in.

Coding agents are the reverse. A code builder is exact, light and straight, but it shows only the parts the agent thought to write. So each object kind is routed one way or the other, and the rule for it changed three times in the case study (Section 5.1). The rule that held is a check, not a list of shapes: a kind may be built in code only when its code build, rendered from in front, shows every part its clean close-up shows, each at least 10 percent visible, with no printed label over a part. Every other kind goes through the prop pipeline: a close-up picture, Pixal3D \citep{li2026pixal3d} in a cloud batch, the paper-thin output closed into a solid, its triangles set by its size, and the straightness check.

\subsection{One model per object}\label{one-model-per-object}

A tool board with twenty tools on it, generated as one model, melts the tools into the board; a console generated as one piece fuses desk, screens, keyboards and cables. SCORE therefore splits composites before anything is generated: one model per real-world object. A parent row (the board, the desk) carries child rows anchored on it (each tool, each monitor), and each child is routed and made on its own. In the central room the tool board became a code-built board with pegs and an outline printed under each tool from that tool's own silhouette, carrying 26 tools.

\subsection{Surfaces: one library, one room wear}\label{surfaces-one-library-one-room-wear}

Generated models carry their own textures, and each brings its own idea of what worn steel looks like; side by side they read as a mix. SCORE keeps shape and surface apart. Models and code give a piece its shape only, and every part names one surface from a shared library. Each surface is a ProcFunc function \citep{raistrick2026procfunc} built on Infinigen's shaders \citep{raistrick2023infinigen}, coloured only from the place's palette tokens, with wear, dirt and seed as named settings, baked in the cloud to the maps the engine needs at the texture density each piece needs. A room has one wear setting, and wear comes from a cause: edges, feet near the floor, drips under pipes. Labels and notes go on top as a few decals, each on a clear flat spot. On a generated model, the surface of each part is chosen from its picture by PartCrafter parts \citep{lin2025partcrafter} laid onto the Pixal3D shape: the picture's colour picks which library surface a whole part takes, never the colour of a single face.

The library for world 1 holds 71 variants in 12 families (painted, steel, aluminium, rubber, plastic, cable, fabric, composite, glass, screen, light and print) from 14 recipes, grown later with Earth surfaces such as plaster, tile, terrazzo, floorboards, rusted steel and wet asphalt.

\subsection{Sound}\label{sound}

Sound is part of the world, not an afterthought, and it is compiled the same way. The world's data says what everything sounds like in three layers: each surface carries its footstep, impact and scrape; each room works out its echo from its size and surfaces; each object's inventory row names its own sound. Every sound is generated in one cloud batch by MOSS-SoundEffect v2.0 \citep{moss2026}, four takes each, and the take is chosen automatically by its match to its prompt under LAION's CLAP model \citep{wu2022clap} less measured faults (clipping, unsteady loops, bad joins). Every take is levelled to one loudness per kind under a true peak of -3 dBTP before it can play. A creator may swap any sound on a review page; picking is optional.

\subsection{Cloud batches and cost control}\label{cloud-batches-and-cost-control}

Every heavy step runs in a rented cloud batch, never on the creator's machine: picture-to-3D, part splitting, surface bakes, sound and, by the end of world 1, the engine's own tests, shots and frame-time bench. Machines are rented for a batch and deleted after it. Each batch records its cost in a ledger, and each place's run records its wall time, its spend and the faults caught before and after spending. A step that grows past its expected size stops itself: when closing a generated piece into a solid threw spikes 60 m long, the step gained a check that stops the job once the copy's box grows more than 3 percent past the model's.

\subsection{Human picks}\label{human-picks}

The creator decides at a few fixed points and nowhere else: the references, the concept, the style, and the final review of each place. Where two routes compete, the review is blind: both sides are drawn the same way (same renderer, light, camera and size), labelled X and Y at random, and the key is kept out of the page until the pick is made. A page that cannot be drawn the same way says that it is not blind. Agents may make interim picks to keep a run moving, but each such pick is marked on its page as the agent's, so the creator can reverse it.

\subsection{The optional world step}\label{the-optional-world-step}

A concept is one picture from one angle. It shows two or three walls, and each object from one side. An optional world step turns the picked concept into a walkable reference of the whole room, from which close-ups of every object can be taken from any side and the walls the concept does not show can be filled in the same style. In our runs this step used World Labs Marble. Nothing a world step makes is shipped: it is a reference for later stages, never part of the world. Where its outputs or anything derived from them are shown, they carry the credit ``Generated using World Labs''.

The world-1 results in Section 5 were all produced \textbf{without} this step: after the route change of 7 October, close-ups were cropped from the concept picture directly. Whether the route makes fuller rooms with the world step than without it is \textcolor{red}{[TODO: not measured; an internal run on one room was started on 8 October and stopped when the game work paused]}.
